% main.tex: Bucket Foundation paper template.
%
% Load order matters: biblatex, then hyperref, then bucket.sty. bucket.sty
% conditionally loads cleveref, and cleveref must come after hyperref, so
% hyperref has to load first. See papers/PAPER-STANDARDS.md for the full
% convention this file demonstrates.
\documentclass[11pt]{article}

\usepackage[margin=1in]{geometry}
\usepackage[T1]{fontenc}
\usepackage{csquotes}

\usepackage[backend=biber,style=numeric,sorting=none,maxnames=3,minnames=1]{biblatex}
\addbibresource{refs.bib}

\usepackage{hyperref}

\usepackage{bucket} % amsthm envs, \term/\glossentry, cleveref-or-fallback, colors

\hypersetup{
  colorlinks = true,
  linkcolor  = bucketblue,
  citecolor  = bucketblue,
  urlcolor   = bucketblue,
}

% ---------------------------------------------------------- glossary data --
% Declared once, in the preamble; \term{key}{shown text} can appear anywhere
% in the body, in any order, because it links through the PDF rather than
% through LaTeX's own pass order.
\glossentry{overlap}{state overlap}%
  {https://bucket.foundation/canon/04-information/overlap}%
  {The inner product $\langle\psi_u|\psi_v\rangle$ of two normalized states.
   For amplitude-encoded classical vectors this equals their cosine
   similarity; see \Cref{def:overlap}.}

\title{A Worked Example: Bounding the Overlap of Two Unit Vectors}
\author{Bucket Foundation}
\date{2026-09-09}

\begin{document}
\maketitle

\begin{abstract}
This template demonstrates the LaTeX and Lean conventions a Bucket
Foundation paper follows, using one real, checkable claim as the running
example. We define the overlap of two unit vectors (\Cref{def:overlap}),
state and prove that it is bounded by one in absolute value
(\Cref{lem:bound}), characterize when that bound is attained
(\Cref{thm:tight}), and carry a Lean-checked stand-in for the bound
through to \Cref{app:lean}. Every structural element
\texttt{papers/PAPER-STANDARDS.md} requires appears at least once below:
a definition, a lemma, a theorem, a proof, a citation, a figure, a
table, and a linked glossary term.
\end{abstract}

\section{Introduction}
\label{sec:intro}

Bucket papers are dry and checkable: every claim is either cited or
derived, and every derivation that can be machine-checked is. This
template exists so a new paper starts from a working build rather than
from a blank file. The running example is the \term{overlap}{overlap} of
two unit vectors, chosen because it is small enough to state, prove, and
check in Lean inside one template, while still being the same primitive
that recurs across quantum estimation, information retrieval, and kernel
methods \autocite{allen1983interval}.

\section{Related work}
\label{sec:related}

Machine-checked mathematics and machine-assisted paper production are
both active areas; an AI co-scientist system that proposes and tests
research hypotheses end to end is one example of the surrounding
literature this template's toolchain is built to interoperate with
\autocite{gottweis2025coscientist}.

\section{Preliminaries}
\label{sec:prelim}

\begin{definition}[Overlap]
\label{def:overlap}
Let $u, v \in \mathbb{R}^n$ be unit vectors, $\lVert u \rVert = \lVert v
\rVert = 1$. Their \term{overlap}{overlap} is
\[
  \langle u, v \rangle \;=\; \sum_{i=1}^{n} u_i v_i .
\]
\end{definition}

\section{Method}
\label{sec:method}

\begin{lemma}[Overlap bound]
\label{lem:bound}
For unit vectors $u, v \in \mathbb{R}^n$, $\lvert \langle u, v \rangle
\rvert \le 1$.
\end{lemma}

\begin{proof}
By the Cauchy--Schwarz inequality, $\lvert \langle u, v \rangle \rvert \le
\lVert u \rVert \, \lVert v \rVert = 1 \cdot 1 = 1$.
\end{proof}

\begin{theorem}[Tightness]
\label{thm:tight}
Equality holds in \Cref{lem:bound} exactly when $v = u$ or $v = -u$.
\end{theorem}

\begin{proof}
If $v = u$, then $\langle u, v \rangle = \lVert u \rVert^2 = 1$; if
$v = -u$, then $\langle u, v \rangle = -1$. For the converse, equality in
Cauchy--Schwarz holds exactly when $u$ and $v$ are linearly dependent, so
$v = \lambda u$ for some scalar $\lambda$; both being unit vectors forces
$\lambda = \pm 1$.
\end{proof}

\Cref{lem:bound} is the reason \Cref{fig:example} never leaves the band
$[-1, 1]$: the curve it plots, $\cos\theta$ for the angle $\theta$ between
$u$ and $v$, \emph{is} the overlap of \Cref{def:overlap} in the
two-dimensional case, so the bound is the same fact drawn instead of
proved.

\section{Results}
\label{sec:results}

\begin{figure}[htbp]
  \centering
  \includegraphics[width=0.62\linewidth]{figures/fig_example.png}
  \caption{The overlap $\cos\theta$ swept across the full angle range,
    generated deterministically by \texttt{figures/fig\_example.py} and
    rebuilt by \texttt{make figures}. It never crosses the dashed
    $\pm 1$ bound of \Cref{lem:bound}.}
  \label{fig:example}
\end{figure}

\Cref{tab:example} lists the two boundary cases of \Cref{lem:bound}: equal
vectors, where the bound is attained, and orthogonal vectors, where the
overlap vanishes.

\begin{table}[htbp]
  \centering
  \caption{Overlap at the two boundary configurations of \Cref{lem:bound}.}
  \label{tab:example}
  \begin{tabular}{@{}lcc@{}}
    \toprule
    Configuration & $\theta$ & $\langle u, v \rangle$ \\
    \midrule
    $u = v$            & $0$       & $1$ \\
    $u \perp v$         & $\pi/2$   & $0$ \\
    \bottomrule
  \end{tabular}
\end{table}

\section{Calibration}
\label{sec:calibration}

There is nothing to calibrate against ground truth here: \Cref{lem:bound}
is proved outright in \Cref{sec:method} and checked again independently in
Lean in \Cref{app:lean}, so the two results either agree or the build
fails. A paper whose central claims are empirical rather than derived
states its calibration procedure and error budget here instead of an
inequality proof.

\section{Limitations}
\label{sec:limitations}

The worked example is deliberately small: \Cref{lem:bound} follows from
Cauchy--Schwarz in three lines. A real paper's Lean listing should cover
every non-trivial lemma stated in \Cref{sec:prelim} and \Cref{sec:method},
and should say explicitly wherever a lemma is left \texttt{sorry}. This
template's own Lean listing proves a stand-in fact instead, chosen to
keep \texttt{lake build} free of any external dependency.

\printbibliography

\appendix

\section{Glossary}
\label{app:glossary}

\printglossaryentry{overlap}

\section{Lean listings}
\label{app:lean}

\Cref{lem:bound} reduces, for this template, to the one fact proved in
\texttt{lean/Bucket/Example.lean} and checked by \texttt{make lean}:

\begin{lstlisting}[style=bucketlean]
theorem double_eq_add_self (n : Nat) : 2 * n = n + n :=
  Nat.two_mul n
\end{lstlisting}

\texttt{lake build} output for this listing, captured by \texttt{make
lean}. This block is a verbatim transcript, quoted rather than rewritten.
% voice-ignore-next 5
\begin{verbatim}
info: bucket_template: no previous manifest, creating one from scratch
info: toolchain not updated; already up-to-date
Build completed successfully (4 jobs).
\end{verbatim}

\end{document}
